\section{GNN-AdaptGraph}\label{sec:method}

Fig.~\ref{fig:arch} shows the adaptation loop. The controller decides \emph{when} and
\emph{how strongly} to adapt. THREAD\_ADJUST, ACGR and METIS decide \emph{how}.

\begin{figure}[t]
\centering
\begin{tikzpicture}[
  node distance=3.5mm and 4mm,
  box/.style={draw=black!45, rounded corners=2pt, align=center, font=\scriptsize,
              minimum height=8.5mm, inner sep=2.5pt, fill=black!3},
  act/.style={draw=black!45, rounded corners=2pt, align=left, font=\scriptsize, inner sep=3pt, fill=white},
  arr/.style={-{Stealth[length=1.6mm]}, black!70, semithick}]
  \node[box, text width=15mm] (stream) {Window $t$\\$\Delta E_t$, touched};
  \node[box, right=of stream, text width=17mm] (ingest) {Ingest\\(LDG places\\new vertices)};
  \node[box, right=of ingest, text width=19mm] (monitor) {Monitor\\8 partition + 3\\global features};
  \node[box, right=of monitor, text width=19mm, fill=blue!6, draw=blue!50] (ctrl) {GraphSAGE\\controller on the\\comm.\ graph};
  \node[act, below=6mm of ctrl, text width=40mm] (acts) {%
    \NOA{}: keep configuration\\
    \TA{}: $\tau_i \propto W_i$ (no migration)\\
    \RB{}: ACGR local refinement\\
    \RP{}: METIS scratch-remap};
  \node[box, left=10mm of acts, text width=32mm] (exec) {BSP executor: per-partition\\thread pools, $S$ supersteps,\\halo exchange of active ghosts};
  \draw[arr] (stream) -- (ingest);
  \draw[arr] (ingest) -- (monitor);
  \draw[arr] (monitor) -- (ctrl);
  \draw[arr] (ctrl) -- node[right, font=\scriptsize] {$a_t$} (acts);
  \draw[arr] (acts) -- (exec);
  \draw[arr] (exec.west) -| node[pos=0.25, below, font=\scriptsize] {next window} (stream.south);
\end{tikzpicture}
\caption{Adaptation loop of GNN-AdaptGraph. Each window is ingested, monitored and
classified by the controller. The chosen action is applied before the window is executed.}
\label{fig:arch}
\end{figure}

\subsection{Runtime State}
The state is the \emph{communication graph} of the current partitioning: one node per
partition, and an edge between two partitions that share cut edges. Each node $i$ carries
eight scale-free features:
\begin{enumerate}
  \item load share $W_i/\bar W$;
  \item thread pressure $(W_i/\tau_i)/\overline{(W/\tau)}$;
  \item ghost ratio $g_i/|P_i|$;
  \item cut ratio $c_i/e_i$;
  \item cut drift $c_i/e_i-c^0_i$ since the last (re)partitioning;
  \item load volatility $|W_i/\bar W-(W_i/\bar W)_{t-1}|$;
  \item size share $|P_i|k/|V|$;
  \item active fraction $|P_i\cap A_t|/|P_i|$.
\end{enumerate}
Here $g_i$ is the number of ghosts held by $P_i$, $c_i$ the number of its cut edges and $e_i$
the number of its incident edges. (The unnormalised features of our first version did not
generalise across graph sizes.) Three global features are broadcast to every node: the work
scale $\log_{10}(1+\sum_i W_i)/7$, the fraction of the window spent communicating, and
$\log_{10}(1+10^3 T_{\mathrm{window}})/4$. They tell the controller how much time is at stake
and whether the platform is compute-bound or communication-bound.

\subsection{GraphSAGE Controller}
The controller $f_\theta$ applies two GraphSAGE layers \cite{hamilton2017graphsage} with the
mean aggregator,
\begin{equation}
  h^{(l)}_{\mathcal{N}(i)} = \tfrac{1}{|\mathcal{N}(i)|}\textstyle\sum_{j\in\mathcal{N}(i)} h^{(l-1)}_j, \qquad
  h^{(l)}_i = \sigma\big(W^{(l)}[\,h^{(l-1)}_i \,\Vert\, h^{(l)}_{\mathcal{N}(i)}\,]\big),
\end{equation}
over the communication graph (input dimension 11, hidden dimension 64, batch normalisation,
a residual connection between the layers). A mean$\Vert$max readout pools the partition
embeddings into a graph embedding $z_t$, and an MLP ($128\!\to\!32\!\to\!4$, dropout 0.3) outputs
one logit per action: $a_t=\arg\max_a f_\theta(z_t)_a$. The model has 14{,}244 parameters.
As an ablation we replace the GraphSAGE layers with per-partition linear layers (MLP
controller, 9{,}444 parameters). This isolates the contribution of neighbourhood aggregation.

\subsection{Outcome-Based Labels and Cost-Sensitive Training}\label{sec:labels}
Our earlier controller was trained to imitate a hand-written threshold rule, so it could not
learn anything the rule did not already encode. The revised controller learns from the
\emph{outcomes} of actions. For a state at window $t$, a look-ahead oracle evaluates every
action $a\in\mathcal{A}$:
\begin{equation}\label{eq:oracle}
  J_H(a) = T^{\mathrm{adapt}}(a)+T^{\mathrm{mig}}(a) + \sum_{h=0}^{H-1} \hat T_{t+h}\big(\text{$a$ at $t$, \NOA{} afterwards}\big),
\end{equation}
where $\hat T$ is the predicted window time (execution plus network) and $H=8$. The label is
$a^\star_t=\arg\min_a J_H(a)$. Because the graph stream does not depend on the actions, the
look-ahead replays exactly the windows that will follow. $\hat T$ comes from a cost model
calibrated on the real executor (Sect.~\ref{sec:executor}):
\begin{equation}\label{eq:cost}
  \hat T^{\mathrm{exec}} = S\big(c_0 + c_1 \max_i W_i/s(\tau_i) + c_2\,\mathrm{CV}_A\big),
\end{equation}
where $s(\tau)$ is the measured parallel speedup with $\tau$ threads. Every sample stores
$J_H(a)$ for all four actions, so we train with the expected regret of the predicted action
distribution, which is the quantity the system cares about:
\begin{equation}\label{eq:loss}
  \mathcal{L}(\theta) = \mathbb{E}_{s}\Big[\textstyle\sum_{a} \pi_\theta(a\,|\,s)\,\big(J_H(a)/J_H(a^\star)-1\big)\Big]
  + 0.1\,\mathrm{CE}\big(\pi_\theta(\cdot\,|\,s), a^\star\big).
\end{equation}
Cross-entropy treats all errors alike, although their costs differ by orders of magnitude.
Triggering METIS when \NOA{} would have been best costs far more than the reverse mistake.
Loss~\eqref{eq:loss} optimises the decision cost directly. In our data, class-weighted
cross-entropy reaches a similar test regret ($\CEregret\pm\CEregretStd$\% vs.\
$\GSregret\pm\GSregretStd$\%, Sect.~\ref{sec:results}). Training states come from an $\varepsilon$-greedy oracle behaviour policy
($\varepsilon=0.3$), so that states off the oracle's path are also covered.

\subsection{ACGR: Communication-Aware Adaptive Graph Repartitioning}\label{sec:acgr}
ACGR repairs an existing partitioning locally by migrating boundary vertices. It minimises
\begin{equation}\label{eq:acgr}
  C = \alpha\,\frac{\mathrm{cut}}{|E|} + \beta\,\frac{\mathrm{CV}}{2|E|}
    + \gamma\,\max\!\big(0,\ \lambda-\epsilon\big) + \delta\,\frac{\min(\bar s,10)}{10},
  \qquad \lambda=\frac{\max_i W_i}{\bar W}-1,
\end{equation}
where $\bar s$ is the mean absolute degree skewness of the partitions and
$(\alpha,\beta,\gamma,\delta)=(0.45,0.35,0.15,0.05)$. Every term is a cost normalised to $[0,1]$,
so lower is better. Our first version gained only one edge on Cora, and two changes fix this.
First, the communication term now uses CV~\eqref{eq:cv}. The old boundary-edge count is
exactly twice the cut, so it merely duplicated the $\alpha$ term. Second, imbalance within the
tolerance $\epsilon$ is a \emph{constraint} rather than a cost, as in METIS. Penalising every
imbalance made each move away from a balanced start look worse. A move $v: P_s\to P_t$ is
therefore admissible only if
\begin{equation}\label{eq:constraint}
  W_t + w_v \le (1+\epsilon)\,\bar W \quad\text{or}\quad W_t+w_v < W_s ,
\end{equation}
and it is accepted only if it strictly decreases $C$.

ACGR keeps, for every vertex, the number of its neighbours in each partition. A move's
change in cut and CV then costs $O(\deg v)$ to evaluate. Its change in imbalance and skew costs
$O(k)$, because the loads and the first three degree moments are kept per partition. Each pass
sorts the boundary vertices by the objective change of their best admissible move. It then
re-evaluates each move against the current state and applies it if it still improves $C$.
A pass thus costs $O(|E|+|V|k)$. Passes repeat until no move improves $C$, up to 3 in the
runtime and 10 in E1. The kernels are compiled with Numba \cite{lam2015numba}. Tests check
the incremental objective against a full recomputation on 50 random graphs.

\subsection{Executor and Real Thread Adjustment}\label{sec:executor}
Each superstep runs Eq.~\eqref{eq:layer} on the active rows of every partition. A row block
is split into $\tau_i$ chunks of equal work and executed by Numba kernels that release the
Python GIL, on a thread pool dedicated to partition $i$. \TA{} therefore really changes the
parallelism of a partition: the thread pools are rebuilt with the new $\tau$ (\ReconfMs\,ms).
Exchanges copy only the values of active boundary vertices into the ghost buffers of the
receiving partitions. The executor's output matches a sequential SciPy reference to
$10^{-5}$ for any partitioning, thread allocation and active set.
